% =========================================================================
\section{Conclusiones}
\label{sec:conclusiones}
% =========================================================================

Puede considerarse que los objetivos planteados en la introducción se han 
cumplido: se ha presentado el formalismo de Mie y de la T-matrix para la
dispersión por una y varias esferas de oro (secciones~\ref{sec:formalismo} y \ref{sec:fundamento multiple}), y se ha
analizado cómo el código MSTM implementa este formalismo y qué
magnitudes físicas calcula (sección~\ref{sec:mstm}). También se ha
desarrollado un entorno propio en Python que automatiza por completo el ciclo
de trabajo con MSTM, y cuyo 
desarrollo ha constituido una parte sustancial del
esfuerzo dedicado al trabajo (sección~\ref{sec:python_wrapper}). Con esta
herramienta se ha estudiado la dispersión por la esfera individual, el dímero
y el tetrámero, analizando tanto el campo lejano como el campo cercano, y,
finalmente, los resultados se han reinterpretado desde el punto de vista de
las nanoantenas ópticas (sección~\ref{sec:nanoantenas}).

Entre los resultados obtenidos destacan varias tendencias. La esfera aislada
produce una intensificación de campo cercano moderada, del orden de diez
veces la intensidad incidente, al carecer de un gap donde concentrar el
campo. La aparición de esa región al pasar al dímero altera significativamente 
la respuesta óptica del sistema: estrechar el hueco aumenta el acoplamiento 
de campo cercano y, con él, la intensificación, 
que en el dímero de $85\,\text{nm}$ alcanza
$\sim\!169\,|\mathbf{E}_0|^2$ para $d=5\,\text{nm}$, acompañada de un
desplazamiento al rojo de la resonancia de más de $170\,\text{nm}$. Este
corrimiento espectral, en cambio, depende del tamaño y resulta prácticamente
nulo en el régimen dipolar ($a=25\,\text{nm}$) y en el régimen Mie avanzado
($a=600\,\text{nm}$). Añadir partículas a la cadena refuerza la concentración
del campo: el tetrámero homogéneo de $25\,\text{nm}$ pasa de los
$147\,|\mathbf{E}_0|^2$ del dímero a $431\,|\mathbf{E}_0|^2$ en su gap
central. La mayor intensificación observada en este trabajo,
$566\,|\mathbf{E}_0|^2$, se obtuvo en el tetrámero heterogéneo
$85$--$25$--$25$--$85$. En conjunto, el tamaño, la separación
entre partículas y la geometría del agregado constituyen los 
grados de libertad que determinan la respuesta óptica de estos sistemas.

El entorno gráfico en Python sobre MSTM es una de las principales 
aportaciones de este trabajo. La herramienta resultó eficaz para el propósito que
motivó su desarrollo: automatizar un flujo de trabajo (conversión de
unidades, generación del fichero de entrada, ejecución del binario y
representación de resultados) que, realizado manualmente, 
supone una tarea laboriosa y susceptible de errores.
El entorno actúa como una capa externa que no modifica MSTM 
y el registro de cada simulación en JSON garantiza su reproducibilidad.
Como contrapartida, el programa hereda las limitaciones del propio MSTM y, en su
estado actual, deja sin explotar parte de las opciones que ya contempla
(otros materiales, haz gaussiano o barridos en radio), empleadas solo
parcialmente en este trabajo.

Por último, el trabajo deja abiertas varias líneas de continuación. La más
inmediata es acotar el tamaño a partir del cual la partícula deja de
comportarse como una nanoantena eficiente, es decir, la transición entre el
régimen de resonancia localizada y el régimen multipolar que aquí solo se ha
identificado de forma cualitativa. También sería natural prolongar el estudio
a cadenas más largas, para observar la saturación del desplazamiento al rojo
descrita en la literatura~\cite{Harris2009,Barrow2011}, así como explorar
otras geometrías, otros materiales ya disponibles en el entorno (Ag, Al, Cu)
y la dependencia con la polarización y con el índice de refracción del medio.
Todas ellas pueden abordarse con la herramienta desarrollada, que queda como punto
de partida para quien desee continuar el trabajo.